\subsection{RQ2: To what extent can SLMs correctly identify errors in generated code?}
\label{sec:rq2}

In the table ~\ref{tab:rq2_detection}, the error detection performance of the evaluated models is summarized, with Accuracy measuring whether the complete set of predicted labels matches those delivered by code execution. Furthermore, Macro-F1 assigns equal importance to the four result categories, CR, FR, RR, and SR, which denote recall for the classifications Correct, Functional Error, Runtime Error, and Syntax Error, respectively. MF1(NS) presents Macro-F1 after excluding syntax errors as a sensitivity analysis. Although the table also shows fitted and larger models, this RQ discusses only the SLMs.

Analyzing the SLMs, it is noticeable that the Muse Glimmer 30B presents the best accuracy and macro-F1 results, with 73.1\% and 53.8\%, respectively, while the Gemma 4 31B presents 72.4\% and 48.5\% in the same metrics, being the second-best. However, unlike the best-performing models, the others show varied combinations in each class's aggregate performance. This difference between Accuracy and Macro-F1 indicates an imbalance: even with high performance on the most frequent error classes, performance remains weak on the less frequent ones.

The two Qwen 2.5 Coder configurations illustrate this distinction. The 7B model achieves higher Accuracy than the 32B model (62.5\% vs. 60.8\%), but the 32B model achieves higher Macro-F1 (41.0\% vs. 33.7\%). The difference mainly stems from the minority error categories: Runtime Recall increases from 2.5\% to 11.4\%, while Syntax Recall increases from 17.6\% to 76.5\%. Therefore, when analyzing Accuracy alone, the differences in prediction distributions across error types go unnoticed.

Detection capability depends on each error category. The Functional Error classification is generally well recognized across a wide range of models, including Gemma 3 4B (84.7\%), the IT version (82.0\%), and Ornith 1.5 35B (78.1\%). However, high Functional Recall does not translate into robust detection, as in Gemma 3 4B, whose Accuracy was only 32.9\%, with 22.6\% Correct Recall and 6.3\% Runtime Recall. Therefore, even if the model recognizes most functional failures, it still struggles to distinguish other results.



Runtime Error classification is the most difficult category for models to detect, with the highest accuracy percentages being 32.9\% for GPT-OSS 20B and 31.6\% for Gemma 4, while the others did not exceed 30\%. This low distribution contrasts sharply with the classification of Functional Errors, for which models are more likely to detect errors. Therefore, Runtime failures remain a problem in error investigation by SLMs.

The Syntax Error classification differs from the others because it shows much more pronounced extremes in detection. For example, GLM 4.7 Flash achieved 94.1\% Syntax Recall; in contrast, Gemma 4 E2B IT and 3 4B IT showed only 5.9\%. Therefore, the ability to detect syntax errors varies widely across models.

The difference between code generation and error detection is quite evident when analyzing GPT-OSS. This is due to the variability in the accuracy of each task, since in RQ1 it achieves 79.0\% Accuracy; however, in detection, the Accuracy is 25.1\%, and the Correct Recall is 17.5\%. Nevertheless, it has the highest Runtime Recall among the models and a Syntax Recall of 72.5\%, despite the low accuracy. Therefore, the low accuracy does not reflect uniformly poor error recognition, but rather asymmetrical recognition.

Since the Syntax Error classification is rare and can be directly identified through deterministic compilation, the Macro-F1 parameter is also presented without this category. This exclusion allows us to assess how this class affects the aggregate macro metric. The largest change was for Muse Glimmer, which rose from 53.8\% to 61.83\%, but others also increased by more than 5\%. Therefore, it is clear that, although rare, Syntax Error can significantly affect Macro-F1.

In addition to detecting errors in code generated by other models, RQ2 also evaluates self-assessment performance. Table \ref{tab:self_detection} presents the main parameter, Self-Det, which measures the proportion of incorrect solutions a model produced but detected itself. The other parameters are the same as those discussed earlier, but in this case they are viewed from the perspective of self-analysis. ND refers to the absence of errors in the evaluated category.

\begin{table}[htbp]
\centering
\caption{Self-error detection by generator model. Self-Det. reports exact-set detection over the model's own erroneous solutions. Class-specific cells report detected/reference counts. ND indicates that the model produced no reference error of that category.}
\label{tab:self_detection}
\begingroup
\setlength{\tabcolsep}{3pt}
\begin{tabular}{lrrrr}
\toprule
\textbf{Model} &
\textbf{Self-Det.} &
\textbf{Funct.} &
\textbf{Runt.} &
\textbf{Synt.} \\
\midrule

\shortstack[l]{Qwen 3.8\\27B}
& \scorecell{67.9}{\shortstack{67.9\%\\(55/81)}}
& \scorecell{75.4}{\shortstack{75.4\%\\(52/69)}}
& \scorecell{25.0}{\shortstack{25.0\%\\(3/12)}}
& ND
\\

\shortstack[l]{Gemma 3\\4B}
& \scorecell{60.4}{\shortstack{60.4\%\\(87/144)}}
& \scorecell{77.3}{\shortstack{77.3\%\\(102/132)}}
& \scorecell{0.0}{\shortstack{0.0\%\\(0/30)}}
& ND
\\

\shortstack[l]{GLM 4.7\\Flash 30B}
& \scorecell{55.0}{\shortstack{55.0\%\\(61/111)}}
& \scorecell{65.6}{\shortstack{65.6\%\\(63/96)}}
& \scorecell{20.0}{\shortstack{20.0\%\\(3/15)}}
& \scorecell{100.0}{\shortstack{100.0\%\\(3/3)}}
\\

\shortstack[l]{Qwen 2.5\\Coder 7B}
& \scorecell{54.4}{\shortstack{54.4\%\\(49/90)}}
& \scorecell{66.7}{\shortstack{66.7\%\\(46/69)}}
& \scorecell{14.3}{\shortstack{14.3\%\\(3/21)}}
& ND
\\

\shortstack[l]{Ornith 1.5\\35B}
& \scorecell{53.8}{\shortstack{53.8\%\\(63/117)}}
& \scorecell{80.0}{\shortstack{80.0\%\\(48/60)}}
& \scorecell{16.7}{\shortstack{16.7\%\\(4/24)}}
& \scorecell{38.5}{\shortstack{38.5\%\\(15/39)}}
\\

\shortstack[l]{Devstral\\Small 2 24B}
& \scorecell{52.3}{\shortstack{52.3\%\\(58/111)}}
& \scorecell{64.4}{\shortstack{64.4\%\\(58/90)}}
& \scorecell{9.1}{\shortstack{9.1\%\\(3/33)}}
& ND
\\

\shortstack[l]{Qwen 3\\Coder 30B}
& \scorecell{42.3}{\shortstack{42.3\%\\(33/78)}}
& \scorecell{51.5}{\shortstack{51.5\%\\(34/66)}}
& \scorecell{0.0}{\shortstack{0.0\%\\(0/15)}}
& ND
\\

\shortstack[l]{Qwen 2.5\\Coder 32B}
& \scorecell{33.3}{\shortstack{33.3\%\\(37/111)}}
& \scorecell{36.9}{\shortstack{36.9\%\\(31/84)}}
& \scorecell{19.4}{\shortstack{19.4\%\\(7/36)}}
& ND
\\

\shortstack[l]{Muse Glimmer\\30B}
& \scorecell{19.0}{\shortstack{19.0\%\\(12/63)}}
& \scorecell{25.9}{\shortstack{25.9\%\\(14/54)}}
& \scorecell{25.0}{\shortstack{25.0\%\\(3/12)}}
& ND
\\

\shortstack[l]{Gemma 4\\31B}
& \scorecell{13.6}{\shortstack{13.6\%\\(9/66)}}
& \scorecell{6.7}{\shortstack{6.7\%\\(3/45)}}
& \scorecell{16.7}{\shortstack{16.7\%\\(3/18)}}
& \scorecell{33.3}{\shortstack{33.3\%\\(3/9)}}
\\

\shortstack[l]{GPT-OSS\\20B}
& \scorecell{5.3}{\shortstack{5.3\%\\(3/57)}}
& \scorecell{15.7}{\shortstack{15.7\%\\(8/51)}}
& \scorecell{0.0}{\shortstack{0.0\%\\(0/21)}}
& ND
\\

\bottomrule
\end{tabular}
\endgroup
\end{table}

Self-error detection varies considerably across the models evaluated, with Qwen 3.8 and GPT-OSS showing accuracies of 67.9\% and 5.3\%, respectively. This variation helps explain the phenomenon and highlights that a good generator is not always a good evaluator, especially of itself: in RQ1, GPT-OSS was one of the best models, with 79.0\% accuracy; however, from this self-analysis perspective, it performed the worst.

The class-specific self-evaluation results follow a pattern similar to general detection. Functional Errors are comparatively easier to recognize, with Ornith reaching 80.0\%, Gemma 3 4B reaching 77.3\%, and Qwen 3.8 reaching 75.4\%. Runtime self-detection remains low: Qwen 3.8 and Muse Glimmer each detect 3 of 12 Runtime cases, GLM detects 3 of 15, and Qwen 2.5 Coder 32B detects 7 of 36, while some models detect none. These percentages must be interpreted together with their denominators because the number of errors produced by each model differs substantially. GLM, for example, recognizes all of its own Syntax Errors, but this corresponds to only three cases.

Overall, the evaluated SLMs exhibit heterogeneous and strongly class-dependent error-detection capabilities. Functional Errors can be recognized with relatively high recall by several models, whereas Runtime Errors remain difficult across the evaluated configurations. Syntax detection varies substantially between models and can noticeably affect Macro-F1 despite the small size of the class. The self-evaluation results further show that code-generation ability and error-recognition ability are distinct: a model that frequently generates correct code may still be unable to reliably identify the errors present in the smaller set of solutions that it gets wrong.
